\documentclass{article}
\usepackage{graphicx} % Required for inserting images
\usepackage{amsmath}
\title{Lesson 6 (Extra Challenge) - Proof Writing}
\author{William Wang}
\date{September 2025}

\begin{document}

\maketitle

\section{}
We have 
\[\sin x = \frac{x^2}{625\pi^2}\]
We know what the graph of the sine function looks like. The graph of the left side should be a massive upwards-opening parabola since it is just a quadratic with $\frac{1}{625\pi^2}$ as the coefficient. The vertex of the parabola is \((0,0)\) so we know that we just want to solve the equation
\[\frac{x^2}{625\pi^2} \le 1\]
since the maximum value of $\sin x$ is 1. 
\[\frac{x^2}{625\pi^2} \le 1 \implies x^2 \le 625\pi^2 \implies x^2-625\pi^2 \le 0 \implies (x-25\pi)(x+25\pi) \le 0\]
Hence
\[x \in [-25\pi, 25\pi]\]
The amount of peaks in this interval is 
\[\text{From } [-25\pi, 0) = -25\pi/2\pi \text{(since the cycle of sine is } 2\pi) = 12\]
\[\text{From } (0, 25\pi] = \mid -25\pi/2\pi \mid \text{(since the cycle of sine is } 2\pi) = 12\]
Additionally, since we skipped 0 and started counting on the positive part of the x axis at $2\pi$, we missed 1 cycle of sine from $(0, 2\pi)$, hence another 2 intersection points. Now we just need to check the two endpoints to see if there is just 1 intersection for that peak. 
\[\sin(-25\pi) \neq 1 \text{ and } \sin(25\pi) \neq 1\]
Hence, the amount of intersections are
\[24\cdot2 +2 = 50\]
\section{}
Skipped. 
\end{document}
